\chapter{Resiliencia, crisis y respaldo externo}
\label{chap:anexo-moderadores}

Este anexo reúne el detalle de los moderadores que la Sección~\ref{sec:moderadores-resumen} resume: la resiliencia de la política económica, el vector de crisis, quién tiene la deuda y el respaldo externo. La deuda, el único moderador con evidencia firme, está en la Sección~\ref{sec:hallazgo-central}.

\section{Resiliencia de la política económica}
\label{sec:robustez-resiliencia}

Esta sección y la siguiente examinan si alguna condición de la economía, además de la deuda (Sección~\ref{sec:hallazgo-central}), modifica la transmisión de cada riesgo al spread. La primera es la resiliencia de la política económica.

La Tabla~\ref{tab:resiliencia} estima la ecuación~\eqref{eq:resiliencia}. Las economías menos resilientes tienen, en promedio, bancos más frágiles y una cola más adversa (correlación del índice con $\ln JLoss_{t-1}$ de $-0{,}43$ y con $D_{t-1}$ de $-0{,}33$).

\begin{table}[htbp]
\centering
\small
\caption[Resiliencia de la política económica]{Resiliencia de la política económica como moderador (efectos fijos de país y de tiempo, $N=670$)}
\label{tab:resiliencia}
\begin{adjustbox}{max width=\linewidth}
\begin{tabular}{l cc}
\toprule
 & Log-log & Niveles (pb) \\
\midrule
Efecto directo de la resiliencia ($\kappa$) & $-0{,}153^{***}$ & $-3{,}5$ \\
Fragilidad $\times$ resiliencia ($\psi_1$) & $0{,}053$ & $4{,}70^{***}$ \\
Riesgo de cola $\times$ resiliencia ($\psi_2$) & $-0{,}014^{*}$ & $-1{,}88$ \\
\quad $p$ \textit{wild cluster bootstrap} ($\psi_1$ / $\psi_2$) & $0{,}42$ / $0{,}46$ & $0{,}22$ / $0{,}43$ \\
\addlinespace
\multicolumn{3}{l}{\textit{Riesgo de cola $\times$ resiliencia, según el régimen}} \\
\quad Fuera de crisis & $-0{,}024^{***}$ & $-4{,}25^{***}$ \\
\quad Cambio en crisis & $0{,}023^{***}$ & $5{,}39^{***}$ \\
\quad En crisis (neto) & $-0{,}001$ & $1{,}14^{*}$ \\
\addlinespace
\multicolumn{3}{l}{\textit{Efecto directo de la resiliencia, según el régimen}} \\
\quad Fuera de crisis & $-0{,}141^{***}$ & $-5{,}6$ \\
\quad En crisis & $-0{,}330^{***}$ & $-25{,}8^{***}$ \\
\addlinespace
Riesgo de cola $\times$ resiliencia, con $R$ en $t-4$ & $-0{,}028^{***}$ & $-4{,}22^{**}$ \\
\bottomrule
\end{tabular}
\end{adjustbox}
\par\smallskip\parbox{0.95\linewidth}{\footnotesize \textit{Nota:} $R$ es el promedio de los valores estandarizados de $-$deuda/PIB, balance fiscal/PIB y reservas/PIB, reestandarizado; los coeficientes son por desviación estándar de resiliencia. Crisis: 2008Q4--2009Q4 y 2020Q1--2021Q4. Niveles: spread en puntos básicos y $JLoss$ en niveles. Errores de Driscoll--Kraay; \textit{wild cluster bootstrap} con pesos de Webb y 13 \textit{clusters}. ***, ** y * indican significancia al 1\%, 5\% y 10\%. Fuente: \texttt{p19\_resiliencia.py}.}
\end{table}

Tres resultados, ninguno más que sugerente. Primero, las economías más resilientes tienen spreads más bajos: una desviación estándar adicional de resiliencia se asocia con un spread en torno a un 14\% menor ($-0{,}153$), más en las crisis. La asociación es significativa solo con Driscoll--Kraay; con errores agrupados por país no lo es ($p=0{,}38$). Segundo, el coeficiente del riesgo de cola es menor donde hay más resiliencia fuera de las crisis ($-0{,}024$ por desviación estándar, una magnitud comparable a la propia semielasticidad media de $0{,}032$), y esa diferencia desaparece durante la crisis financiera global y la pandemia ($-0{,}001$). Una lectura posible es la de un seguro: un mejor resguardo abarata la prima mientras el riesgo es idiosincrásico, pero cuando el choque golpea a todos a la vez las primas suben para todos. Es una lectura, no un hallazgo: el contraste por régimen es significativo solo con Driscoll--Kraay. Tercero, para la fragilidad bancaria el signo es el contrario al esperado: su coeficiente es mayor, no menor, en las economías más resilientes ($+0{,}053$, sin significancia; $+4{,}7$ puntos básicos en niveles, solo con Driscoll--Kraay). Ninguno de los términos de moderación del índice resiste el \textit{wild bootstrap} sobre la muestra completa ($p=0{,}42$ y $0{,}46$ en logaritmos; $0{,}22$ y $0{,}43$ en puntos básicos).

\subsection{Qué componente de la resiliencia importa}
\label{sec:resiliencia-variantes}

El índice $R$ tiene dos debilidades: promedia tres variables que no tienen por qué operar igual, y se mide en el mismo período que el spread, de modo que puede estar respondiendo a él. La Tabla~\ref{tab:resiliencia-variantes} lo abre de tres formas: cada componente por separado, su primer componente principal, y una clasificación \emph{ex ante} que separa a las economías en alta y baja resiliencia según su promedio de 2004--2007 y no cambia en el tiempo.

\begin{table}[htbp]
\centering
\small
\caption[Componentes de la resiliencia]{Moderación por componente de la resiliencia (efectos fijos de país y de tiempo, $N=670$; $p$ de \textit{wild cluster bootstrap} entre paréntesis)}
\label{tab:resiliencia-variantes}
\begin{adjustbox}{max width=\linewidth}
\begin{tabular}{l cc cc}
\toprule
 & \multicolumn{2}{c}{Log-log} & \multicolumn{2}{c}{Niveles (pb)} \\
\cmidrule(lr){2-3}\cmidrule(lr){4-5}
Moderador & Fragilidad ($\psi_1$) & Riesgo de cola ($\psi_2$) & Fragilidad ($\psi_1$) & Riesgo de cola ($\psi_2$) \\
\midrule
Índice $R$ (Tabla~\ref{tab:resiliencia}) & $0{,}053$ ($0{,}42$) & $-0{,}014^{*}$ ($0{,}46$) & $4{,}70^{***}$ ($0{,}22$) & $-1{,}88$ ($0{,}43$) \\
\addlinespace
Reservas/PIB & $-0{,}247^{***}$ ($0{,}08$) & $-0{,}001$ ($0{,}97$) & $-11{,}33^{***}$ ($0{,}04$) & $0{,}05$ ($0{,}98$) \\
Menor deuda/PIB & $0{,}107^{***}$ ($0{,}21$) & $-0{,}010$ ($0{,}06$) & $4{,}55^{***}$ ($0{,}12$) & $-1{,}51$ ($0{,}17$) \\
Balance fiscal/PIB & $0{,}048$ ($0{,}31$) & $0{,}003$ ($0{,}36$) & $3{,}81^{***}$ ($0{,}15$) & $0{,}34$ ($0{,}71$) \\
\addlinespace
Primer componente principal & $0{,}094^{***}$ ($0{,}22$) & $-0{,}008$ ($0{,}32$) & $4{,}50^{***}$ ($0{,}15$) & $-1{,}19$ ($0{,}38$) \\
\addlinespace
Alta resiliencia \textit{ex ante} & $-0{,}022$ ($0{,}89$) & $-0{,}028^{***}$ ($0{,}06$) & $-7{,}86^{***}$ ($0{,}27$) & $-4{,}73^{**}$ ($0{,}15$) \\
\quad sin trimestres de crisis ($N=576$) & $-0{,}051$ ($0{,}71$) & $-0{,}047^{***}$ ($0{,}08$) & $-7{,}16^{***}$ ($0{,}24$) & $-8{,}52^{***}$ ($0{,}11$) \\
\bottomrule
\end{tabular}
\end{adjustbox}
\par\smallskip\parbox{0.95\linewidth}{\footnotesize \textit{Nota:} Cada fila reemplaza $R$ en la ecuación~\eqref{eq:resiliencia} por el moderador indicado, estandarizado y con signo tal que un valor mayor es más resiliencia. El primer componente principal explica el 58\% de la varianza conjunta, con cargas de $0{,}69$ (deuda), $0{,}70$ (balance fiscal) y $0{,}18$ (reservas). La clasificación \textit{ex ante} es un indicador de país (resiliencia sobre la mediana en el promedio de 2004--2007 para las ocho economías con datos, y en sus primeros dieciséis trimestres para el resto); su efecto directo lo absorben los efectos fijos de país. Alta resiliencia: Chile, China, Indonesia, Malasia, Perú y Sudáfrica. Estrellas según Driscoll--Kraay: ***, ** y * al 1\%, 5\% y 10\%. \textit{Wild cluster bootstrap} con pesos de Webb y 1\,999 réplicas. Fuente: \texttt{p19\_resiliencia.py}.}
\end{table}

La Tabla~\ref{tab:resiliencia-variantes} contiene doce términos de moderación por forma funcional, y dos destacan: las reservas con la fragilidad, y la resiliencia previa con el riesgo de cola. Ninguno de los dos estaba preespecificado, y por eso se someten a un escrutinio adicional antes de leerlos.

\textbf{Reservas y fragilidad bancaria.} El coeficiente de la fragilidad es menor donde hay más reservas: $-0{,}247$ por desviación estándar de reservas sobre PIB en logaritmos ($p$ de \textit{wild bootstrap} de $0{,}08$) y $-11{,}3$ puntos básicos en niveles ($0{,}04$). La Tabla~\ref{tab:reservas} lo examina por cinco lados.

\begin{table}[htbp]
\centering
\small
\caption[Reservas y fragilidad bancaria bajo escrutinio]{El término reservas $\times$ fragilidad bajo escrutinio: rezagos, exclusión de Polonia, versión \textit{ex ante}, comparaciones múltiples y moderadores placebo}
\label{tab:reservas}
\begin{adjustbox}{max width=\linewidth}
\begin{tabular}{l cccc ccc}
\toprule
 & \multicolumn{4}{c}{Log-log} & \multicolumn{3}{c}{Niveles (pb)} \\
\cmidrule(lr){2-5}\cmidrule(lr){6-8}
Moderador $\times$ fragilidad & Coef. & $p$ DK & $p$ \textit{cluster} & $p$ \textit{wild} & Coef. & $p$ \textit{cluster} & $p$ \textit{wild} \\
\midrule
\multicolumn{8}{l}{\textit{A. Reservas/PIB}} \\
\quad Contemporáneas ($N=670$) & $-0{,}247$ & $<0{,}001$ & $0{,}048$ & $0{,}08$ & $-11{,}3$ & $<0{,}001$ & $0{,}04$ \\
\quad Rezagadas cuatro trimestres ($N=681$) & $-0{,}234$ & $<0{,}001$ & $0{,}057$ & $0{,}051$ & $-10{,}9$ & $<0{,}001$ & $0{,}02$ \\
\quad Rezagadas ocho trimestres ($N=673$) & $-0{,}190$ & $<0{,}001$ & $0{,}12$ & $0{,}085$ & $-10{,}3$ & $0{,}001$ & $0{,}057$ \\
\quad Contemporáneas, sin Polonia ($N=583$) & $-0{,}084$ & $<0{,}001$ & $0{,}19$ & $0{,}30$ & $-9{,}5$ & $<0{,}001$ & $0{,}08$ \\
\quad \textit{Ex ante}, invariantes en el tiempo ($N=765$) & $-0{,}073$ & $0{,}035$ & $0{,}16$ & $0{,}30$ & $-4{,}3$ & $<0{,}001$ & $0{,}13$ \\
\quad Exclusiones de un país con $p$ \textit{wild} $<0{,}05$ & \multicolumn{4}{c}{1 de 13} & \multicolumn{3}{c}{10 de 13} \\
\quad $p$ \textit{wild} con corrección de Holm (seis contrastes) & \multicolumn{4}{c}{$0{,}40$} & \multicolumn{3}{c}{$0{,}22$} \\
\addlinespace
\multicolumn{8}{l}{\textit{B. Moderadores placebo}} \\
\quad Cuenta corriente/PIB, contemporánea & $-0{,}087$ & $0{,}049$ & $0{,}046$ & $0{,}084$ & $-5{,}3$ & $<0{,}001$ & $0{,}042$ \\
\quad Cuenta corriente/PIB, rezagada cuatro trimestres & $-0{,}113$ & $0{,}011$ & $0{,}012$ & $0{,}040$ & $-4{,}1$ & $0{,}007$ & $0{,}062$ \\
\quad Tipo de cambio real, contemporáneo & $+0{,}039$ & $0{,}048$ & $0{,}37$ & $0{,}33$ & $+0{,}1$ & $0{,}94$ & $0{,}95$ \\
\quad Tipo de cambio real, rezagado cuatro trimestres & $+0{,}123$ & $<0{,}001$ & $0{,}024$ & $0{,}019$ & $+3{,}6$ & $0{,}002$ & $0{,}026$ \\
\quad Placebos significativos al 5\%, de nueve: DK / \textit{cluster} / \textit{wild} & \multicolumn{4}{c}{7 / 3 / 2} & \multicolumn{3}{c}{6 / 4 / 2} \\
\bottomrule
\end{tabular}
\end{adjustbox}
\par\smallskip\parbox{0.95\linewidth}{\footnotesize \textit{Nota:} Cada fila reemplaza $R$ en la ecuación~\eqref{eq:resiliencia} por el moderador indicado, estandarizado; el coeficiente es el de su producto con la fragilidad. \textit{Ex ante}: media del país en 2004--2007 (ocho economías) o en sus primeros dieciséis trimestres (cinco); su efecto directo lo absorben los efectos fijos de país. Holm: familia de tres componentes de la resiliencia por dos riesgos. Los nueve placebos son la cuenta corriente, la inflación, el tipo de cambio real efectivo y la concentración bancaria trimestral, contemporáneos y rezagados cuatro trimestres, y la concentración bancaria estructural. \textit{Cluster}: errores agrupados por país. \textit{Wild cluster bootstrap} por país con pesos de Webb. Fuente: \texttt{p31\_reservas\_placebo.py}.}
\end{table}

A su favor: la asociación es estable al rezagar las reservas cuatro y ocho trimestres ($-0{,}234$ y $-0{,}190$, con $p$ de \textit{wild bootstrap} de $0{,}051$ y $0{,}085$), de modo que no proviene de que las reservas respondan al spread del mismo período; y en puntos básicos no depende de Polonia ($-9{,}5$; $p=0{,}08$; significativa en diez de las trece exclusiones de un país). En su contra hay cuatro resultados. En logaritmos depende de Polonia: sin ella el coeficiente cae a un tercio ($-0{,}084$) y no es significativo con errores agrupados por país ($p=0{,}19$) ni con \textit{wild bootstrap} ($0{,}30$), y solo una de las trece exclusiones pasa el \textit{wild bootstrap} al 5\%. La versión \textit{ex ante}, con las reservas fijadas antes de la muestra, cae a $-0{,}073$ ($p=0{,}30$). Con la corrección de \citet{Holm1979} sobre la familia examinada ---tres componentes por dos riesgos--- el $p$ es $0{,}40$ en logaritmos y $0{,}22$ en puntos básicos. Y los placebos: la cuenta corriente rezagada ($-0{,}113$; $p=0{,}040$) y el tipo de cambio real rezagado ($+0{,}123$; $p=0{,}019$; en puntos básicos, $0{,}026$) muestran interacciones con la fragilidad de fuerza similar o mayor, y con errores de Driscoll--Kraay entre seis y siete de nueve moderadores arbitrarios resultan significativos. La forma lineal tampoco es creíble en todo el soporte (Figura~\ref{fig:reservas-efmarg}): la elasticidad implícita de la fragilidad cruza el cero en $+0{,}08$ desviaciones estándar de reservas y es significativamente negativa para el 22,5\% de las observaciones en logaritmos y el 28,8\% en puntos básicos.


En suma: el coeficiente de la fragilidad es menor donde hay más reservas; la asociación es estable al rezagar las reservas y se mantiene en puntos básicos, donde es más firme que en logaritmos y no depende de Polonia; pero en logaritmos depende de Polonia, no pasa la corrección por comparaciones múltiples en ninguna de las dos formas, otras variables macroeconómicas muestran interacciones de fuerza similar y la forma lineal implica una elasticidad negativa que no es creíble. Se trata como exploratoria, y no sostiene ninguna afirmación sobre lo que ocurriría si una economía acumulara reservas.

\textbf{Resiliencia previa y riesgo de cola.} En las economías clasificadas como de alta resiliencia \textit{ex ante}, el coeficiente del riesgo de cola es $0{,}028$ menor ($p$ de \textit{wild bootstrap} de $0{,}064$) y $4{,}7$ puntos básicos menor en niveles ($0{,}15$). El signo es estable ---negativo en las trece exclusiones de un país---, pero sin Polonia pierde significancia en logaritmos ($-0{,}018$; $p=0{,}11$ con Driscoll--Kraay) y en puntos básicos desaparece sin México. La clasificación, además, no es anterior a la muestra para cinco de las trece economías (Brasil, Chile, Colombia, México y Perú usan sus primeros dieciséis trimestres, 2009Q4--2013Q3). Es también exploratorio.

\textbf{Espacio fiscal.} Con menos deuda, el coeficiente de la fragilidad es mayor y no menor ($+0{,}107$), que es lo que hacía positivo el coeficiente del índice: el primer componente principal carga casi solo en las dos variables fiscales. No resiste el \textit{wild bootstrap} ($p=0{,}21$).

La resiliencia no modifica, entonces, la transmisión de ninguno de los dos riesgos de forma demostrable. Lo que queda es sugerente (el spread más bajo de las economías resilientes y el contraste por régimen del riesgo de cola, solo con Driscoll--Kraay) o exploratorio (las reservas y la resiliencia previa).

\section{Crisis, respaldo externo y quién tiene la deuda}
\label{sec:crisis-interaccion}

Las Tablas~\ref{tab:crisis-lnlag-embiext} y~\ref{tab:crisis-lnlag-embiext-b}, en el Anexo~\ref{chap:anexo-robustez}, estiman la ecuación~\eqref{eq:crisisint} en cuatro modelos anidados (CM1: $\ln JLoss$ y su interacción con crisis; CM2: lo mismo con $D$; CM3: ambos niveles y sus interacciones; CM4: la especificación completa) bajo las tres estructuras de efectos fijos, con el vector de crisis único (Tabla~\ref{tab:crisis-lnlag-embiext}) y descompuesto en \textit{Backstop} y \textit{EMstress} (Tabla~\ref{tab:crisis-lnlag-embiext-b}).

Fuera de crisis, la interacción es indistinguible de cero ($\hat\beta_3$ entre $-0{,}002$ y $+0{,}013$ en CM4). Bajo \textit{Backstop} ---crisis financiera global y pandemia--- la suma $\hat\beta_3+\hat\beta_4$ es negativa ($-0{,}015$, $p=0{,}03$ con efectos fijos de país y de tiempo) y bajo \textit{EMstress} ---2015--2016--- positiva ($+0{,}026$, $p=0{,}08$). Una lectura de esa asimetría es que el respaldo oficial de la Reserva Federal y del FMI desacopla la transmisión doméstica. Cuatro pruebas indican que el panel no la sostiene.

\textbf{Primero, errores agrupados por país.} Con errores \textit{cluster} por país en lugar de Driscoll--Kraay, ninguna de las dos sumas es significativa: $-0{,}015$ con $p=0{,}40$ bajo \textit{Backstop} y $+0{,}026$ con $p=0{,}11$ bajo \textit{EMstress}.

\textbf{Segundo, otros episodios sin respaldo.} Si el mecanismo opera en el estrés emergente sin respaldo oficial, debería aparecer también en los otros episodios de ese tipo en la muestra: el \textit{taper tantrum} (2013Q2--Q3) y la venta masiva de activos emergentes de 2018 (2018Q2--Q3). No aparece: la suma es $-0{,}020$ ($p=0{,}68$) en el primero y $+0{,}037$ ($p=0{,}42$) en el segundo, y con los tres episodios juntos es $+0{,}013$ ($p=0{,}25$).

\textbf{Tercero, permutación.} Reestimando la especificación con cada una de las 51 ventanas de tres trimestres consecutivos fuera de \textit{Backstop} (con al menos cinco economías en la muestra) en lugar del episodio de 2015--2016, la suma observada ocupa el lugar 20: una de cada tres ventanas elegidas al azar produce una amplificación igual o mayor ($p=0{,}39$; Figura~\ref{fig:permutacion}). El estrés emergente de 2015--2016 no es un episodio especial en esta dimensión.


\textbf{Cuarto, comparaciones múltiples.} La interacción se contrasta en varias muestras y regímenes. Con la corrección de \citet{Holm1979} sobre esa familia de seis contrastes, la suma bajo \textit{Backstop} pasa de $p=0{,}03$ a $0{,}12$ y la de \textit{EMstress} de $0{,}08$ a $0{,}24$ (Sección~\ref{sec:inferencia-potencia}).

La Figura~\ref{fig:crisisregimen} traza la elasticidad de la fragilidad en función de $D$ en cada régimen: las bandas se superponen en todo el soporte. Las asimetrías por régimen de crisis no son, por tanto, un hallazgo. Las tres subsecciones siguientes examinan las dos explicaciones que se habían propuesto para ellas ---quién tiene la deuda y el respaldo externo--- con variables que las miden directamente.


\subsection{Heterogeneidad por tipo de economía: una agrupación \textit{ex post}}
\label{sec:robustez-heterogeneidad}

La Tabla~\ref{tab:heterogeneidad} estima la ecuación~\eqref{eq:heterogeneidad}: la interacción en un grupo de once economías y en el grupo formado por Polonia e India, en ambas formas funcionales. La agrupación es \textit{ex post}: proviene de una versión anterior del análisis, en la que la interacción se diluía al añadir esas dos economías.

\begin{table}[htbp]
\centering
\small
\caption[Heterogeneidad por tipo de economía: agrupación \textit{ex post}]{Interacción en una agrupación \textit{ex post}: un grupo de once economías, y Polonia e India (M4, efectos fijos de país y de tiempo, sin controles salvo indicación)}
\label{tab:heterogeneidad}
\begin{adjustbox}{max width=\linewidth}
\begin{tabular}{l cc cc}
\toprule
 & \multicolumn{2}{c}{Log-log ($\beta_3$)} & \multicolumn{2}{c}{Niveles ($\gamma_3$, pb)} \\
\cmidrule(lr){2-3}\cmidrule(lr){4-5}
 & Completa & Sin crisis & Completa & Sin crisis \\
\midrule
Grupo de once economías & $0{,}011$ & $0{,}034^{***}$ & $0{,}78^{**}$ & $1{,}79^{***}$ \\
\quad $p$ \textit{wild cluster bootstrap} & $0{,}12$ & $0{,}015$ & $0{,}046$ & $0{,}002$ \\
\quad con seis controles domésticos & $0{,}004$ & $0{,}004$ & $0{,}53^{***}$ & $0{,}81^{***}$ \\
\quad $p$ \textit{wild} con controles & $0{,}39$ & $0{,}68$ & $0{,}042$ & $0{,}12$ \\
\quad \textit{leave-one-country-out} (rango) & $[0{,}002;\,0{,}015]$ & $[0{,}018;\,0{,}039]$ & $[0{,}27;\,0{,}93]$ & $[0{,}87;\,1{,}92]$ \\
\addlinespace
Polonia e India (modelo con interacciones de grupo) & $-0{,}059^{**}$ & $-0{,}118^{***}$ & $0{,}02$ & $-1{,}82^{**}$ \\
\addlinespace
Observaciones (grupo de once) & 621 & 533 & 621 & 533 \\
\bottomrule
\end{tabular}
\end{adjustbox}
\par\smallskip\parbox{0.95\linewidth}{\footnotesize \textit{Nota:} Errores de Driscoll--Kraay salvo indicación; \textit{wild cluster bootstrap} restringido con pesos de Webb y 1\,999 réplicas (11 \textit{clusters}). Sin crisis: excluye 2008Q4--2009Q4 y 2020Q1--2021Q4. La fila de Polonia e India es $\beta_3+\theta_3$ en la ecuación~\eqref{eq:heterogeneidad}, estimada sobre las trece economías. ***, ** y * indican significancia al 1\%, 5\% y 10\%. Fuente: \texttt{p16\_heterogeneidad.py}, \texttt{paper\_heterogeneidad\_numeros*.csv}.}
\end{table}

Fuera de los trimestres de crisis, en el grupo de once economías la interacción es positiva en ambas formas funcionales ---$+0{,}034$ en elasticidades y $+1{,}8$ puntos básicos en niveles---, resiste el \textit{wild bootstrap} ($p=0{,}015$ y $0{,}002$), se mantiene positiva al excluir cada economía del grupo y pasa la corrección de Holm sobre los contrastes por muestra y régimen ($p$ corregido de $0{,}005$ con Driscoll--Kraay). En Polonia e India es negativa. Es el resultado numéricamente más favorable a la hipótesis de amplificación, y se reporta completo. Pero se rotula exploratorio, por cuatro razones. Primera, la división entre grupos se definió después de ver los datos, de modo que ni la significancia ni la corrección de Holm ---que no cuenta las agrupaciones que pudieron probarse--- valen lo que indican. Segunda, en logaritmos los seis controles domésticos la absorben por completo ($0{,}004$); en niveles se mantiene con ellos, pero con controles y sin crisis ya no resiste el \textit{wild bootstrap} ($p=0{,}12$). Tercera, sobre la muestra completa la interacción en logaritmos no es significativa en ese grupo ($0{,}011$; $p=0{,}12$). Y cuarta, la más importante: la variable que debía explicar la agrupación ---quién tiene la deuda--- no lo hace, como muestra la subsección siguiente.

\subsection{Participación extranjera en la deuda: una prueba preespecificada}
\label{sec:participacion-extranjera}

La interpretación que acompañaba a la agrupación anterior era que el pasivo contingente bancario se traslada al spread donde el inversionista marginal de la deuda es un extranjero sensible al riesgo. Esa interpretación se puede contrastar con un moderador continuo y medido de forma independiente: la participación de no residentes en la deuda del gobierno, de la base de \citet{ArslanalpTsuda2014} (versión de mayo de 2026; trece países, 2004Q1--2025Q4). El moderador mide la tenencia extranjera de la deuda total del gobierno (en moneda local y extranjera), mientras que el EMBI es el spread de la deuda externa en dólares; la prueba sería más directa con la participación extranjera en la deuda en moneda local, que la base cubre de forma incompleta.

\textbf{Preespecificación.} Las hipótesis se escribieron el 13 de octubre de 2026, antes de estimar y de ver ningún resultado, y no se agregaron pruebas después. La preespecificación consta en el encabezado del propio código (\texttt{p26\_participacion\_extranjera.py}), no en un registro externo fechado. El moderador primario, $F_{i,t-4}$, es la participación de los tenedores extranjeros no bancarios en la deuda total del gobierno general, rezagada cuatro trimestres y estandarizada (media de 21,5\% y desviación estándar de 14,4 puntos porcentuales; el 75\% de su varianza es entre países). La especificación es
\begin{multline}
\ln\mathrm{EMBI}_{i,t}=\alpha_i+\delta_t+\beta_1\ln JLoss_{i,t-1}+\beta_2D_{i,t-1}+\beta_3(\ln JLoss\times D)_{i,t-1}+f\,F_{i,t-4}\\
+h_1(\ln JLoss_{i,t-1}\times F_{i,t-4})+h_2(D_{i,t-1}\times F_{i,t-4})+h_3\big((\ln JLoss\times D)_{i,t-1}\times F_{i,t-4}\big)+\varepsilon_{i,t}.
\label{eq:participacion}
\end{multline}
La familia primaria son los tres coeficientes $h_1$, $h_2$ y $h_3$, con signo predicho positivo, corrección de Holm sobre los tres y veredicto por \textit{wild cluster bootstrap} por país (pesos de Webb, 1\,999 réplicas). Las variantes ---participación extranjera total, versión \textit{ex ante}, rezago de ocho trimestres, muestra sin crisis y descomposición entre y dentro de países--- se declararon secundarias. Y se fijó de antemano la consecuencia: si $h_3$ no resultaba significativo en la prueba primaria, la heterogeneidad de la Sección~\ref{sec:robustez-heterogeneidad} pasaba a exploratoria.

\begin{table}[htbp]
\centering
\small
\caption[Participación extranjera en la deuda como moderador]{Participación extranjera en la deuda del gobierno como moderador (efectos fijos de país y de tiempo)}
\label{tab:participacion}
\begin{adjustbox}{max width=\linewidth}
\begin{tabular}{l ccc}
\toprule
 & Fragilidad $\times F$ ($h_1$) & Riesgo de cola $\times F$ ($h_2$) & Interacción $\times F$ ($h_3$) \\
\midrule
\multicolumn{4}{l}{\textit{A. Familia primaria: no bancos extranjeros, rezago de cuatro trimestres ($N=762$)}} \\
\quad Log-log & $+0{,}037$ & $-0{,}001$ & $-0{,}006$ \\
\quad \quad $p$: Driscoll--Kraay / \textit{cluster} / \textit{wild} & $0{,}34$ / $0{,}39$ / $0{,}41$ & $0{,}90$ / $0{,}91$ / $0{,}92$ & $0{,}31$ / $0{,}49$ / $0{,}48$ \\
\quad \quad $p$ \textit{wild} con corrección de Holm & $1{,}00$ & $1{,}00$ & $1{,}00$ \\
\quad Niveles (pb) & $+3{,}38$ & $+1{,}51$ & $-0{,}33$ \\
\quad \quad $p$: Driscoll--Kraay / \textit{cluster} / \textit{wild} & $0{,}002$ / $0{,}013$ / $0{,}13$ & $0{,}12$ / $0{,}32$ / $0{,}34$ & $0{,}053$ / $0{,}31$ / $0{,}38$ \\
\quad \quad $p$ \textit{wild} con corrección de Holm & $0{,}40$ & $0{,}69$ & $0{,}69$ \\
\addlinespace
\multicolumn{4}{l}{\textit{B. Secundarias, log-log ($p$ de \textit{wild bootstrap} entre paréntesis)}} \\
\quad Participación extranjera total & $+0{,}035$ ($0{,}54$) & $-0{,}005$ ($0{,}41$) & $-0{,}006$ ($0{,}58$) \\
\quad \textit{Ex ante} (media del país en 2004--2007) & $-0{,}017$ ($0{,}76$) & $+0{,}013$ ($0{,}13$) & $-0{,}020$ ($0{,}038$) \\
\quad Rezago de ocho trimestres & $+0{,}042$ ($0{,}30$) & $+0{,}002$ ($0{,}75$) & $-0{,}013$ ($0{,}19$) \\
\quad Sin trimestres de crisis ($N=653$) & $+0{,}069$ ($0{,}17$) & $+0{,}010$ ($0{,}13$) & $-0{,}003$ ($0{,}87$) \\
\quad Solo variación entre países & $+0{,}027$ ($0{,}73$) & $-0{,}000$ ($0{,}98$) & $-0{,}014$ ($0{,}085$) \\
\quad Solo variación dentro del país & $+0{,}049$ ($0{,}38$) & $+0{,}004$ ($0{,}76$) & $+0{,}028$ ($0{,}30$) \\
\bottomrule
\end{tabular}
\end{adjustbox}
\par\smallskip\parbox{0.95\linewidth}{\footnotesize \textit{Nota:} Ecuación~\eqref{eq:participacion}; coeficientes por desviación estándar del moderador. Signo predicho: positivo en las tres columnas. Datos de \citet{ArslanalpTsuda2014}, versión de mayo de 2026. El rezago de cuatro trimestres hace perder tres observaciones. \textit{Cluster}: errores agrupados por país. \textit{Wild cluster bootstrap} por país con pesos de Webb y 1\,999 réplicas. No se ponen estrellas: el veredicto preespecificado es el del \textit{wild bootstrap} con corrección de Holm. Fuente: \texttt{p26\_participacion\_extranjera.py}.}
\end{table}

\textbf{Resultado.} Según el criterio declarado ---el \textit{wild bootstrap} con corrección de Holm---, ninguna de las tres hipótesis se sostiene (Tabla~\ref{tab:participacion}). En logaritmos, $h_1=+0{,}037$ ($p$ de \textit{wild bootstrap} de $0{,}41$), $h_2=-0{,}001$ ($0{,}92$) y $h_3=-0{,}006$ ($0{,}48$), este último con el signo contrario al predicho; con la corrección de Holm los tres $p$ son $1{,}00$. En puntos básicos el término con la fragilidad es $+3{,}38$, significativo con Driscoll--Kraay ($p=0{,}002$) y con errores agrupados por país ($0{,}013$) pero no con el \textit{wild bootstrap} ($0{,}13$; con Holm, $0{,}40$): por la escala de este capítulo es sugerente en puntos básicos, y nulo según el criterio preespecificado. Los otros dos términos no son significativos ($0{,}34$ y $0{,}38$). Las pruebas secundarias no cambian el cuadro: con la participación total los coeficientes son los mismos; la versión \textit{ex ante} da un $h_3$ de $-0{,}020$ ($p=0{,}038$), significativo y de signo contrario al predicho; y con el rezago de ocho trimestres o sin los trimestres de crisis nada resiste el \textit{wild bootstrap}.

\textbf{Consecuencia.} Tal como se declaró de antemano, la heterogeneidad del grupo de once economías pasa a exploratoria. El resultado de esta prueba es ausencia de evidencia con intervalos anchos, no evidencia de un cero: como tres cuartas partes de la varianza del moderador son entre países, la prueba compara en lo esencial trece economías.



\textbf{Un hecho descriptivo que corrige la etiqueta.} La Figura~\ref{fig:participacion} y la Tabla~\ref{tab:participacion-ranking} ordenan las economías por la participación de no residentes en su deuda. La participación total media es de 64\% en Perú, 50\% en Indonesia, 43\% en Colombia, 40\% en Polonia, 39\% en México, 31\% en Turquía, Chile y Sudáfrica, 29\% en Filipinas, 24\% en Malasia, 12\% en Brasil, 5\% en India y 3\% en China. Polonia es la cuarta de trece (con un rango de 27\% a 55\% en la muestra), y la tercera en no bancos, con 29,5\%. Una versión anterior de este trabajo describía a Polonia e India como economías cuya deuda se coloca en el mercado local; la descripción vale para India, no para Polonia. Y China y Brasil, que están en el grupo de once, son de las de menor participación extranjera. La agrupación de la Sección~\ref{sec:robustez-heterogeneidad} no corresponde, entonces, a quién tiene la deuda: es una agrupación \textit{ex post}, y así se la llama en este trabajo.

\subsection{Respaldo externo medido por país}
\label{sec:respaldo-pais}

La lectura del vector de crisis descansa en una variable que solo cambia en el tiempo: \textit{Backstop} vale uno para todas las economías en la crisis financiera global y en la pandemia, aunque solo Brasil y México tuvieron una línea \textit{swap} con la Reserva Federal. Con ella no se puede separar el respaldo externo de cualquier otra cosa que haya ocurrido en esos trimestres. El contraste creíble usa variación entre países: qué economías tenían, en cada trimestre, acceso efectivo a liquidez oficial en dólares. Se construyó un indicador $Acc_{i,t}$ que vale uno si el país tenía vigente una línea \textit{swap} con la Reserva Federal (Brasil y México, 2008Q4--2009Q4 y 2020Q1--2021Q4) o una Línea de Crédito Flexible del FMI (México desde 2009Q2; Colombia, 2009Q2--2025Q1; Polonia, 2009Q2--2017Q3; Chile desde 2020Q2; Perú, 2020Q2--2024Q2). Las fechas provienen de los comunicados de la Reserva Federal y del FMI y, para las líneas de 2020, de la base del Programa de Estabilidad Financiera de la Universidad de Yale. El indicador no incluye la facilidad de recompra FIMA de la Reserva Federal (desde marzo de 2020) ni las líneas bilaterales con otros bancos centrales. Se estimó
\begin{multline}
\ln\mathrm{EMBI}_{i,t}=\alpha_i+\delta_t+\beta_1\ln JLoss_{i,t-1}+\beta_2D_{i,t-1}+\beta_3(\ln JLoss\times D)_{i,t-1}+a\,Acc_{i,t}\\
+c_1(\ln JLoss_{i,t-1}\times Acc_{i,t})+c_2(D_{i,t-1}\times Acc_{i,t})+c_3\big((\ln JLoss\times D)_{i,t-1}\times Acc_{i,t}\big)+\varepsilon_{i,t}.
\label{eq:respaldo}
\end{multline}
Si el respaldo desacopla la transmisión, cada riesgo debería pesar menos con acceso ($c_1<0$, $c_2<0$) y la interacción ser menor ($c_3<0$).

\begin{table}[htbp]
\centering
\small
\caption[Respaldo externo medido por país]{Respaldo externo por país y fecha como moderador (efectos fijos de país y de tiempo; $p$ de Driscoll--Kraay y de \textit{wild cluster bootstrap} entre paréntesis)}
\label{tab:respaldo}
\begin{adjustbox}{max width=\linewidth}
\begin{tabular}{l cccc}
\toprule
 & Acceso ($a$) & Fragilidad $\times$ acceso ($c_1$) & Cola $\times$ acceso ($c_2$) & Interacción $\times$ acceso ($c_3$) \\
\midrule
\multicolumn{5}{l}{\textit{Log-log}} \\
\quad Muestra completa ($N=765$; 223 con acceso) & $0{,}196^{*}$ & $0{,}127$ ($0{,}14$; $0{,}58$) & $0{,}002$ ($0{,}92$; $0{,}96$) & $-0{,}004$ ($0{,}86$; $0{,}90$) \\
\quad Solo Línea de Crédito Flexible (208) & $0{,}228^{**}$ & $0{,}169^{*}$ & $-0{,}000$ & $-0{,}004$ \\
\quad Dentro de las crisis globales ($N=109$; 54) & $0{,}553^{**}$ & $0{,}063$ ($0{,}78$; $0{,}80$) & $-0{,}018$ ($0{,}21$; $0{,}15$) & $0{,}045$ ($0{,}03$; $0{,}22$) \\
\addlinespace
\multicolumn{5}{l}{\textit{Niveles (pb)}} \\
\quad Muestra completa & $7{,}9$ & $1{,}51$ ($0{,}69$; $0{,}86$) & $6{,}19$ ($0{,}04$; $0{,}27$) & $-1{,}02$ ($0{,}22$; $0{,}52$) \\
\quad Dentro de las crisis globales & $24{,}4$ & $1{,}94$ ($0{,}55$; $0{,}54$) & $-1{,}63$ ($0{,}30$; $0{,}08$) & $0{,}14$ ($0{,}55$; $0{,}79$) \\
\bottomrule
\end{tabular}
\end{adjustbox}
\par\smallskip\parbox{0.95\linewidth}{\footnotesize \textit{Nota:} Seis economías tienen acceso en algún trimestre. Crisis globales: 2008Q4--2009Q4 y 2020Q1--2021Q4. La línea \textit{swap} por sí sola cubre 26 observaciones de dos países y no se reporta por separado. El acceso de Chile es continuo desde 2020Q2: en mayo de 2022 pasó a una Línea de Liquidez de Corto Plazo y el 29 de agosto de 2022 el FMI aprobó una nueva Línea de Crédito Flexible, con la que Chile canceló la anterior. \textit{Wild cluster bootstrap} con pesos de Webb y 1\,999 réplicas. ***, ** y * indican significancia al 1\%, 5\% y 10\% (Driscoll--Kraay). Fuente: \texttt{p25\_respaldo\_externo.py} y \texttt{respaldo\_externo.csv}.}
\end{table}

El resultado no respalda la lectura del respaldo externo. Sobre la muestra completa ninguno de los tres términos tiene el signo y la significancia que predice el desacople: con acceso, el coeficiente de la fragilidad es algo mayor y no menor ($c_1=+0{,}13$, sin significancia), el del riesgo de cola es el mismo y la interacción no cambia ($c_3=-0{,}004$; $p=0{,}86$; \textit{wild bootstrap}, $0{,}90$). Dentro de las crisis globales ---donde la comparación es entre países con y sin acceso en el mismo trimestre, que es el contraste más limpio--- la interacción es \emph{mayor} en los países con acceso ($c_3=+0{,}045$), el signo opuesto al esperado, aunque no resiste el \textit{wild bootstrap} ($p=0{,}22$). El único indicio a favor está en el riesgo de cola, cuyo coeficiente es menor con acceso dentro de las crisis ($c_2<0$), sin alcanzar significancia.

Dos razones limitan la prueba, y conviene declararlas. La primera es que el acceso no es exógeno: el coeficiente directo es positivo (los países con acceso tienen spreads más altos), lo que sugiere que las líneas se piden y se otorgan cuando el riesgo es mayor, o que identifican a economías más integradas y expuestas al ciclo global. La segunda es que son seis economías. Aun así, la conclusión es clara en lo que importa para este trabajo: la asimetría entre \textit{Backstop} y \textit{EMstress} de la Tabla~\ref{tab:crisis-lnlag-embiext-b} no se reproduce cuando el respaldo se mide donde efectivamente estuvo disponible. No hay evidencia de que el respaldo externo desacople la transmisión.

\textbf{Balance de las tres condiciones.} Ni la resiliencia de la política económica, ni el respaldo externo, ni quién tiene la deuda modifican de forma demostrable la transmisión de los dos riesgos al spread. Para la participación extranjera no hay evidencia según el criterio preespecificado (solo un indicio sugerente en puntos básicos para la fragilidad), y para el respaldo externo tampoco; para las reservas, la resiliencia previa y la agrupación \textit{ex post} hay resultados exploratorios que no pasan las pruebas que los harían firmes. En todos los casos se trata de ausencia de evidencia con trece países e intervalos anchos, no de evidencia de que esas condiciones sean irrelevantes.
